\section{The machine and its execution relation}
\label{sec:isa}

\subsection{MIPS32 as a guest}

Throughout this paper MIPS32 means Release 2 of the specification, the revision this work targets and the last before Release~6 re-encoded the instruction set incompatibly. The guest machine is a 32-bit MIPS32 core with 32 integer registers (register 0 hard-wired to zero), the \texttt{HI}/\texttt{LO} pair written by multiplication and division, and a little-endian byte-addressed memory. Three features of the ISA shape the arithmetisation more than anything else.

\paragraph{Branch-delay slots.} The instruction that follows a branch or jump is executed before control transfers. The machine state therefore carries a program counter \emph{pair}: the address of the current instruction and the address of the one that follows it. A sequential instruction advances the pair as $(\mathsf{pc}, \mathsf{next\_pc}) \mapsto (\mathsf{next\_pc}, \mathsf{next\_pc} + 4)$; a branch instruction sets $\mathsf{next\_next\_pc}$ to the target or to $\mathsf{next\_pc}+4$. Consequently the successor of a sequential instruction is \emph{data}, not $\mathsf{pc}+4$: when the instruction sits in a delay slot its successor is the branch target. Every instruction row carries $\mathsf{next\_pc}$ as a witnessed column for this reason.

\paragraph{Unaligned and narrow accesses.} \texttt{LB}, \texttt{LBU}, \texttt{LH}, \texttt{LHU}, \texttt{SB}, \texttt{SH} address individual bytes and half-words of a word, and \texttt{LWL}/\texttt{LWR}/\texttt{SWL}/\texttt{SWR} merge the left or right part of an unaligned word with a register. Memory is modelled at word granularity, so these instructions select and merge bytes of a word in-circuit; they get chips of their own (Section~\ref{sec:air-chips}).

\paragraph{Bit manipulation.} \texttt{EXT}, \texttt{INS}, \texttt{WSBH}, \texttt{SEB}/\texttt{SEH}, \texttt{ROTR}, \texttt{CLO}/\texttt{CLZ} and the conditional moves are common in hashing and serialisation code. All of them are supported natively rather than lowered.

\subsection{Syscalls and precompiles}

Guest programs reach the host through \texttt{SYSCALL}. Besides input/output the syscall table exposes \emph{precompiles}: SHA-256, Keccak-f[1600], the $\Poseidon{}$ permutation, curve arithmetic on secp256k1, secp256r1, BN254, BLS12-381 and Ed25519, field arithmetic for BN254 and BLS12-381, 256-bit modular and $256\times 2048$-bit multiplication. A precompile reads and writes guest memory directly; each invocation becomes one or more rows of a dedicated chip whose memory accesses are timestamped at the invoking cycle.

\paragraph{Hash functions.} Hashing is the dominant primitive of the workloads this class of machine is asked to prove, and it motivates the accelerator interface. An unaccelerated guest computing SHA-256 or Keccak in MIPS instructions pays tens of instructions per round, each occupying an instruction-chip row, whereas an accelerator computes a permutation in a handful of wide rows. Three hash families are accelerated. SHA-256 is split into its message-schedule extension and compression rounds, with a control chip sequencing memory accesses. Keccak-f[1600] is a sponge chip with its own controller and, at 2\,637 columns, is one of the two widest chips (Table~\ref{tab:chips}). $\Poseidon{}$ is accelerated for guests that verify proofs of this system. The same permutation also implements Merkle commitments and the Fiat--Shamir transcript (Section~\ref{sec:prelim}) and appears in the recursion machine (Section~\ref{sec:recursion}). Hint data written by the host into untouched memory becomes part of the memory's initial state (Section~\ref{sec:air-memory}), so it remains subject to the argument's range discipline.

\subsection{Clock, shards and the execution record}

Execution is divided into \emph{shards}. Within a shard a clock $\mathsf{clk}$ advances by five per instruction; a syscall additionally advances it by a bounded number of ticks so that the memory accesses of a precompile fit strictly between the calling instruction and its successor. An instruction executes at $\mathsf{clk}$ and its memory and register accesses occupy the four \emph{sub-cycle positions} $\mathsf{clk}+1,\dots,\mathsf{clk}+4$, one per operand role (Section~\ref{sec:air-frame}), so a $(\mathsf{shard}, \mathsf{clk} + \mathsf{pos})$ pair identifies every access uniquely and the memory argument timestamps an access by $\mathsf{clk}+\mathsf{pos}$. The clock is bounded by $2^{25}$, which caps a shard at $2^{25}/5 \approx 6.71$\,M cycles whatever the configured shard size says. The clock is witnessed as a 16-bit limb and a 10-bit high limb, the same decomposition the memory argument range-checks. A shard therefore closes on whichever of three bounds binds first: that fence; the configured shard size, a few million cycles, against which the executor compares $\mathsf{clk} + \mathsf{max\_syscall\_cycles}$ after scaling the configured figure by four, since it is given in cycles and the clock advances by five per instruction; or a chip's height budget. Which one binds depends on the configuration, and Section~\ref{sec:performance} reports the shard counts that result from four of them.

At a shard boundary the emulator emits, for every memory word touched in the shard, the value and timestamp of its first and last access; these are the leaves of the cross-shard memory argument. The prover arithmetises a shard record containing per-chip event lists, the touched-address ledger and public values. Those public values contain the committed-output and deferred-proof digests, entry and exit program-counter pairs and clocks, exit code, shard indices, global-memory cursors and row counts, and the septic digest of Section~\ref{sec:air-global}. They are the interface through which consecutive shards and recursion layers are chained.

\subsection{The execution relation}
\label{sec:isa-emulator}

Write $\mathsf{Emulate}(P, \mathsf{in}, \mathsf{hint}) = (\mathsf{out}, e)$ for the relation this section defines: the machine of \S\ref{sec:isa} started on the program image $P$ with the two streams available to it, run until it halts, produces the output stream $\mathsf{out}$ and the exit status $e$. It is deterministic: no step reads a clock, a random source or any host state other than the streams.

The relation is realised twice, and the two realisations must agree. The reference is an interpreter: it implements every instruction, every syscall and every error path, and it is the only one that emits the event record the prover consumes. A just-in-time compiler serves the two jobs that dominate wall-clock time and need no events, namely running a guest for its output alone and producing, per shard, the descriptor from which a worker replays that shard. Three checks connect them, none of which is a proof of equivalence for every program: the descriptors of a compiled run are byte-identical to an interpreted one on the tested corpus, the specification vectors of Section~\ref{sec:verification} run on both, and continuous integration compares digests of the emitted records. The difference that matters to the rest of the paper is rate, because a host feeding several accelerators is bound by the executor rather than by the accelerator (\S\ref{sec:perf-host}); Appendix~\ref{sec:executor} gives the mechanism.
